\documentclass[14pt,a4paper,oneside]{extbook}
\usepackage[utf8]{vietnam}
\usepackage{amsmath,amssymb,mathrsfs,fancyhdr,enumerate,multirow,makecell}
\usepackage{tkz-euclide,tikz-3dplot,tikz,tkz-tab}
\usetikzlibrary{arrows.meta,calc,intersections,angles,quotes,backgrounds,shapes.geometric,patterns,through,shadows,positioning}
\usepackage{yhmath}
\usepackage[top=0.7cm, bottom=0.7cm, left=1.0cm, right=1.0cm, headheight=16pt, headsep=3mm, footskip=5mm]{geometry}
\usepackage[hidelinks,unicode]{hyperref}
\usepackage{qrcode}
\usepackage[many]{tcolorbox}
\usepackage{fontawesome}

\pagestyle{fancy}
\fancyhf{}
\fancyhead[L]{\footnotesize\bfseries\color{teal!85!black}\faBook~SỔ TAY KIẾN THỨC TOÁN 4}
\fancyhead[R]{\footnotesize\bfseries\color{blue!85!black}Lớp Toán Thầy Lĩnh -- loptoanthaylinh.vn}
\fancyfoot[C]{\footnotesize\bfseries\color{gray!80!black}-- Sổ tay kiến thức: Bài 46. Tìm số trung bình cộng --}
\renewcommand{\headrulewidth}{0.4pt}

\begin{document}

% --- BANNER TIÊU ĐỀ BÀI HỌC ---
\noindent
\begin{tcolorbox}[
  enhanced,
  colback=teal!8!white,
  colframe=teal!80!black,
  arc=3mm, boxrule=1pt,
  left=3mm, right=3mm, top=2mm, bottom=2mm,
  halign=center
]
  {\Large\bfseries\color{teal!90!black}\faBookmark\ BÀI 46: TÌM SỐ TRUNG BÌNH CỘNG}\\[1.5pt]
  {\footnotesize\itshape\color{black!75} Khái niệm chia đều, công thức thuận, công thức ngược và phương pháp giải toán trên 1 trang A4}
\end{tcolorbox}

\vspace{2mm}

% --- KHỐI 1 & 2: LÝ THUYẾT & KĨ NĂNG ---
\noindent
\begin{minipage}[t]{0.49\textwidth}
  \begin{tcolorbox}[
    enhanced,
    title={\footnotesize\bfseries\color{white}\faBalanceScale\ 1. Bản chất \& Công thức thuận},
    colback=blue!2!white,
    colframe=blue!75!black,
    arc=2.5mm, boxrule=0.8pt, top=1.8mm, bottom=1.8mm, left=2mm, right=2mm
  ]
    \footnotesize
    $\bullet$ \textbf{Bản chất}: Trung bình cộng là giá trị thu được khi lấy \textbf{tổng chia đều} cho tất cả các phần.\\
    $\bullet$ \textbf{Công thức thuận}:
    \begin{center}
      $\mathbf{\text{TBC} = \frac{\text{Tổng các số hạng}}{\text{Số các số hạng}}}$
    \end{center}
    --- Tính tổng các số hạng trước (đóng ngoặc đơn).\\
    --- Đếm xem có bao nhiêu số hạng rồi chia cho số đó.
  \end{tcolorbox}
\end{minipage}%
\hfill
\begin{minipage}[t]{0.49\textwidth}
  \begin{tcolorbox}[
    enhanced,
    title={\footnotesize\bfseries\color{white}\faUndo\ 2. Công thức ngược \& Tính chất},
    colback=teal!2!white,
    colframe=teal!75!black,
    arc=2.5mm, boxrule=0.8pt, top=1.8mm, bottom=1.8mm, left=2mm, right=2mm
  ]
    \footnotesize
    $\bullet$ \textbf{Công thức ngược (cực kì quan trọng)}:
    \begin{center}
      $\mathbf{\text{Tổng} = \text{TBC} \times \text{Số các số hạng}}$
    \end{center}
    Biết TBC của $n$ số $\rightarrow$ tính được ngay Tổng!\\
    $\bullet$ \textbf{Tính chất dãy số cách đều}:\\
    --- Số số hạng lẻ $\rightarrow$ TBC chính là \textbf{số ở chính giữa}.\\
    --- Luôn có: $\text{TBC} = (\text{Số đầu} + \text{Số cuối}) : 2$.
  \end{tcolorbox}
\end{minipage}

\vspace{2mm}

% --- KHỐI 3: SƠ ĐỒ TƯ DUY TIKZ ---
\noindent
\begin{tcolorbox}[
  enhanced,
  title={\footnotesize\bfseries\color{white}\faSitemap\ SƠ ĐỒ TƯ DUY: HAI CHIỀU CÔNG THỨC TRUNG BÌNH CỘNG},
  colback=yellow!2!white,
  colframe=orange!85!black,
  arc=2.5mm, boxrule=0.8pt, top=1.5mm, bottom=1.5mm, left=1mm, right=1mm,
  halign=center
]
  \begin{tikzpicture}[
    scale=0.92,
    cbox/.style={rectangle, rounded corners=2mm, draw=teal!80!black, fill=teal!10, thick, inner sep=3pt, font=\scriptsize\bfseries, align=center, text width=4.8cm},
    arr/.style={-{Stealth[scale=0.85]}, thick, draw=black!70, shorten >= 2.5pt, shorten <= 2.5pt}
  ]
    \node[cbox, draw=blue!75!black, fill=blue!10, text width=4.5cm] (left) at (-0.5, 0) {\textbf{\color{blue!85!black}TỔNG CÁC SỐ HẠNG}\\$(a + b + c + \dots)$\\Cộng tất cả các số};
    \node[cbox, draw=teal!80!black, fill=teal!10, text width=4.5cm] (right) at (9.5, 0) {\textbf{\color{teal!85!black}SỐ TRUNG BÌNH CỘNG}\\Giá trị bình quân chia đều\\$\text{Số bé} < \text{TBC} < \text{Số lớn}$};

    \draw[arr, draw=orange!85!black] ($(left.east)+(0, 0.4)$) -- node[above=2pt, font=\scriptsize\bfseries\color{orange!85!black}]{Chia cho Số các số hạng} ($(right.west)+(0, 0.4)$);
    \draw[arr, draw=purple!85!black] ($(right.west)+(0, -0.4)$) -- node[below=2pt, font=\scriptsize\bfseries\color{purple!85!black}]{Nhân với Số các số hạng} ($(left.east)+(0, -0.4)$);
  \end{tikzpicture}
\end{tcolorbox}

\vspace{2mm}

% --- KHỐI 4: VÍ DỤ MINH HỌA & BẪY TRẮC NGHIỆM ---
\noindent
\begin{minipage}[t]{0.53\textwidth}
  \begin{tcolorbox}[
    enhanced,
    title={\footnotesize\bfseries\color{white}\faCheckCircle\ Ví dụ thực hành chuẩn mực},
    colback=green!2!white,
    colframe=green!60!black,
    arc=2.5mm, boxrule=0.8pt, top=1.8mm, bottom=1.8mm, left=2.2mm, right=2.2mm
  ]
    \footnotesize
    \textbf{Bài toán thuận}: Tìm TBC của $45, 65, 70$:\\
    --- Tổng: $45 + 65 + 70 = 180$.\\
    --- Có $3$ số hạng $\Rightarrow \text{TBC} = 180 : 3 = \mathbf{60}$.\\[1mm]
    \textbf{Bài toán ngược}: Trung bình cộng số tuổi của hai anh em là $12$ tuổi. Anh $15$ tuổi. Tính tuổi em:\\
    --- Tổng số tuổi 2 anh em: $12 \times 2 = 24$ (tuổi).\\
    --- Tuổi của em: $24 - 15 = \mathbf{9}$ (tuổi).
  \end{tcolorbox}
\end{minipage}%
\hfill
\begin{minipage}[t]{0.45\textwidth}
  \begin{tcolorbox}[
    enhanced,
    title={\footnotesize\bfseries\color{white}\faExclamationTriangle\ Mẹo tránh bẫy trắc nghiệm},
    colback=red!2!white,
    colframe=red!70!black,
    arc=2.5mm, boxrule=0.8pt, top=1.8mm, bottom=1.8mm, left=2.2mm, right=2.2mm
  ]
    \footnotesize
    $\bullet$ \textbf{Bẫy thiếu dấu ngoặc đơn}:\\
    Viết biểu thức $45 + 65 + 70 : 3$ là \textbf{SAI}. Bắt buộc phải viết $(45 + 65 + 70) : 3$.\\
    $\bullet$ \textbf{Bẫy đếm sai số các số hạng}:\\
    Luôn đếm cẩn thận có đúng bao nhiêu số hạng trước khi thực hiện phép chia!\\
    $\bullet$ TBC không thể nhỏ hơn số bé nhất hoặc lớn hơn số lớn nhất.
  \end{tcolorbox}
\end{minipage}

\end{document}
